% BHU PYQ -- Physics Mechanics & General Properties of Matter (Mid Term) (2024-25 NEP)
\documentclass[11pt,a4paper]{article}
\usepackage[a4paper,top=2.5cm,bottom=2.5cm,left=2.8cm,right=2.8cm,headheight=14pt]{geometry}
\usepackage{amsmath,amssymb}
\usepackage[T1]{fontenc}
\usepackage[utf8]{inputenc}
\usepackage{lmodern,microtype,fancyhdr,enumitem,hyperref,lastpage,parskip}

\hypersetup{colorlinks=true,urlcolor=black,linkcolor=black,pdftitle=PHYMJ11: Mechanics & General Properties of Matter (Mid Term) -- BHU NEP 2024-25}

\pagestyle{fancy}\fancyhf{}
\renewcommand{\headrulewidth}{0.4pt}\renewcommand{\footrulewidth}{0.4pt}
\fancyhead[L]{\small\textit{Banaras Hindu University}}
\fancyhead[R]{\small\textit{PHYMJ11: Mechanics & General Properties of Matter (Mid Term) (2024-25)}}
\fancyfoot[C]{\small Page \thepage\ of \pageref{LastPage}}

\newlist{parts}{enumerate}{2}
\setlist[parts,1]{label=(\alph*),leftmargin=*,itemsep=4pt,topsep=3pt}
\setlist[parts,2]{label=(\roman*),leftmargin=*,itemsep=2pt,topsep=2pt}
\newcommand{\pts}[1]{\hfill\textit{[#1]}}

\begin{document}
\begin{center}
  {\large\bfseries BANARAS HINDU UNIVERSITY}\\[2pt]
  {\normalsize U.G. Semester 1 (NEP) Examination, 2024-25}\\[6pt]
  \rule{0.8\linewidth}{0.5pt}\\[6pt]
  {\large\bfseries PHYSICS}\\[4pt]
  {\large\bfseries Paper: PHYMJ11 --- Mechanics & General Properties of Matter (Mid Term)}\\[4pt]
  {\normalsize Time: 2:30 Hours \hfill Full Marks: 70}
\end{center}
\rule{\linewidth}{0.4pt}
\smallskip
\noindent\textit{\noindent\textit{Instructions: Attempt \textbf{THREE} questions including Question No.~1, which is compulsory. The figures in the right-hand margin indicate marks.}
\bigskip

\noindent\textbf{Question 1.} \textit{(Compulsory)}
\begin{parts}
  \item State and explain Newton's laws of motion.\pts{3}
  \item Derive the expression for the centre of mass of a system of $N$ discrete particles.\pts{4}
\end{parts}

\bigskip
\noindent\textbf{Question 2.}
\begin{parts}
  \item Define work, energy, and power. State and prove the Work-Energy Theorem.\pts{7}
  \item Show that total linear momentum is conserved in the absence of external forces.\pts{7}
\end{parts}

\bigskip
\noindent\textbf{Question 3.}
\begin{parts}
  \item Distinguish between elastic and inelastic collisions in one and two dimensions.\pts{7}
  \item A particle of mass $m_1$ moving with velocity $v_1$ collides elastically with a stationary particle of mass $m_2$. Calculate the final velocities after collision.\pts{7}
\end{parts}


\vfill\rule{\linewidth}{0.4pt}
\end{document}
